\begin{textsample}{POS Dim 4 – persona\_grok – Score 14.00 – t194\_grok.txt}  \label{ex:f4_pos_040}
The abandoned FSO Safer oil tanker, moored off Yemen 's Red Sea \textbf{coast}, poses dire risks to the \textbf{region}.
This decaying \textbf{vessel} holds over 1.1 million barrels of crude oil, which could lead to devastating \textbf{spills} or explosions.
Such disasters would severely impact human \textbf{life}, \textbf{wildlife}, \textbf{fisheries}, and \textbf{ecosystems}. \textbf{Fishing} communities in the \textbf{area} would face particularly severe effects, with \textbf{livelihoods} threatened by contaminated \textbf{waters} and disrupted marine resources.
Cleanup efforts following a \textbf{spill} would generate massive amounts of hazardous waste, complicating recovery and posing long-term environmental challenges.
The oil itself carries significant health dangers due to its toxic and carcinogenic properties.
Exposure could affect organs, skin, and lungs, leading to serious health issues for those in proximity.
To mitigate risks, individuals should use protective \textbf{equipment} when necessary, coordinate with authorities for any response \textbf{activities}, avoid contaminated \textbf{areas}, and refrain from burning waste or attempting informal cleanups.
Greenpeace welcomes the recent agreement to transfer the oil from the FSO Safer but urges the immediate deployment of containment booms and standby response \textbf{equipment} to prevent catastrophe.
This situation underscores the failures of the oil industry, and we call for its end to avert future \textbf{threats} like this.

% matched lemmas: activity, area, coast, ecosystem, equipment, fishery, fishing, life, livelihood, region, spill, threat, vessel, water, wildlife
\end{textsample}
